\section{The exact identities of the original transport kernels}
\label{sec:kernel-transport-identities}
We now identify the kernels estimated in the preceding sections with
the transports used in the quadratic construction. This also accounts
for the harmonic mode. Throughout this section $k>0$, the block side
$L$ is odd and at least three, and the adjacent train orders $j,t$
satisfy $4\le t\le j\le t+1$.

\subsection{Covariance recursion and exterior algebra}
Fix a nonzero real coarse momentum $p$ in the common real cube
containing the Brillouin cube, and write $q_l=(p+2\pi l)/L$
for its centered aliases. In the following formulas the fine exterior
symbol is rescaled: $\delta_r=Ld_r(q_l)$, whereas $d_r=d_r(p)$
is the coarse symbol. Stars denote Hermitian adjoints at these real
momenta. Let $H_l,H$ be the fine and coarse two-form pulse symbols.
Let $D_l$ be the diagonal matrix with entries $L^{-2}D_{f,a}(q_l)$,
and $D$ the diagonal matrix with entries $D_{c,a}(p)$, where
$D_{v,a}$ is the propagated one-form covariance from
Section~\ref{sec:ex-lattice-free-kernel-components}.
Set
\[
 x=d_0,\quad x_l=\delta_0,\quad
 R=d_1^*Hd_1,\quad R_l=\delta_1^*H_l\delta_1.
\]
All diagonal entries of $D,D_l$ are positive. Put
\[
 \rho=x^*D^{-1}x,\quad
 \pi=\rho^{-1}x^*D^{-1},\quad
 c=\rho^{-1}D^{-1}x.
\]
The corresponding fine objects have subscript $l$. The pulse formula
is $H=\rho^{-1}\bigwedge^2D^{-1}$; with the stated scaling,
$H_l=\rho_l^{-1}\bigwedge^2D_l^{-1}$ as well. Consequently
\begin{equation}\label{eq:transport-rotor}
 R=D^{-1}-\rho^{-1}D^{-1}xx^*D^{-1},\qquad
 DR=I-x\pi,\qquad RD=I-cx^*.
\end{equation}
For completeness, if $D^{-1}=\operatorname{diag}(f_a)$, the
diagonal entry of $d_1^*(\bigwedge^2D^{-1})d_1$ at $a$ is
$f_a\sum_{b\ne a}f_b|x_b|^2$, and its off-diagonal entry at
$(a,b)$ is $-f_af_bx_a\overline{x_b}$. Division by
$\rho=\sum_bf_b|x_b|^2$ proves the first equality; multiplication
by $D$ proves the other two. In particular $Rx=0$ and $x^*R=0$.

Let $m_l$ be the patch mass filter at $q_l$ and define
\[
 F_{r,l}=m_l\bigwedge^r\operatorname{diag}
                 (\operatorname{rail}_L(q_{l,a})).
\]
The geometric-series identity
$L\operatorname{rail}_L(q_{l,a})(e^{iq_{l,a}}-1)=e^{ip_a}-1$
gives
\begin{equation}\label{eq:transport-filter-chain}
 F_{r+1,l}\delta_r=d_rF_{r,l}.
\end{equation}
The covariance recursion is
\begin{equation}\label{eq:transport-covariance-recursion}
                  \sum_l F_{1,l}D_lF_{1,l}^*=D-k^{-1}I.
\end{equation}
Here the sum is a sum over all $L^4$ aliases, with no averaging
factor. To verify the normalization and the recursion, the feed of a
train of total side $M$ is
\[
 M^{-2}\sum_{u\in\Lambda_M}
 |\operatorname{drive}_{v,a}(q_u)|^2 f(q_u).
\]
On appending a block, its drive is the old drive times the new
patch filter evaluated at the old rescaled momentum. Decomposing a
residue modulo $LM$ into its residues modulo $L$ and $M$ therefore
turns this feed into the successive two feeds: the factors are
$(LM)^{-2}=L^{-2}M^{-2}$, and the summands agree by periodicity.
The noise-prefix list of the new train consists of its new unit
prefix and the old prefixes with the block appended. Thus its noise
feed is one plus the patch feed of the old noise. Applying these two
facts to the free covariance and to each earlier independent noise
gives $D_{c,a}=k^{-1}+\operatorname{slide}_aD_{f,a}$.
The diagonal entries of \eqref{eq:transport-covariance-recursion}
are precisely this equality. All off-diagonal entries vanish.

\subsection{The two punctured transport formulas}
Define scalars
\[
 a=\sum_l |m_l|^2\operatorname{tr}D_l,\qquad
 e_l=\overline{m_l}\operatorname{tr}D_l/a,
 \qquad y_l=|m_l|^2\rho/\rho_l.
\]
The principal mass filter is nonzero on the real cube. Since each
$D_l$ is positive, its contribution shows $a>0$. In particular
\begin{equation}\label{eq:transport-weight-sum}
                              \sum_l e_lm_l=1.
\end{equation}
The source constructs the following rescaled columns:
\begin{align}
 \mathcal A_l&=D_lF_{1,l}^*R+e_lx_l\pi,
                                      \label{eq:transport-A}\\
 \mathcal T_l&=H_l\delta_1D_lF_{1,l}^*d_1^*
                    +y_lF_{2,l}^*(I-C),
 & C&=Hd_1Dd_1^*.                    \label{eq:transport-T}
\end{align}
The complement $I-C$ is part of the two-form transport; it cannot
be discarded merely because it vanishes on weighted exact forms.
Indeed \eqref{eq:transport-rotor} and $d_1x=0$ imply
\[
 RDd_1^*=d_1^*,\qquad d_1DR=d_1,
 \qquad d_1^*(I-C)=0,\qquad (I-C)Hd_1=0.
\]
Also, the adjoint of \eqref{eq:transport-filter-chain} gives
$x_l^*F_{1,l}^*d_1^*=0$. These identities prove
\begin{equation}\label{eq:transport-punctured-three}
 R_l\mathcal A_l=F_{1,l}^*R,\qquad
 \delta_1^*\mathcal T_l=F_{1,l}^*d_1^*,\qquad
 \mathcal T_lHd_1=H_l\delta_1\mathcal A_l.
\end{equation}
Here are the cancellations explicitly. The second term in
$R_l\mathcal A_l$ vanishes because $R_lx_l=0$; in the first,
$R_lD_l=I-c_lx_l^*$ and $x_l^*F_{1,l}^*R=0$.
For the second identity the first term is
$R_lD_lF_{1,l}^*d_1^*=F_{1,l}^*d_1^*$, while the second is
$y_lF_{1,l}^*d_1^*(I-C)=0$.
For the third identity the complement term vanishes by
$(I-C)Hd_1=0$, and the remaining term is
$H_l\delta_1D_lF_{1,l}^*R$. The other side equals this
because $\delta_1x_l=0$.

Finally \eqref{eq:transport-filter-chain} at degree zero,
\eqref{eq:transport-weight-sum}, and
\eqref{eq:transport-covariance-recursion} give
\begin{align}
 \sum_l F_{1,l}\mathcal A_l
  &=(D-k^{-1}I)R+x\pi\notag\\
  &=I-k^{-1}R.                       \label{eq:transport-alias-sum}
\end{align}
This is the identity that incorporates the new Gaussian noise.

\subsection{Removing the puncture without changing the kernels}
We check that these formulas are the analytic columns whose bounds
were proved earlier. In this continuation every real-momentum
adjoint is written as transposition after reflection $p\mapsto-p$;
for example, $|m(p)|^2$ continues as $m(p)m(-p)$.
For an alias-indexed factor the reflection is $q_l\mapsto-q_l$,
equivalently $(p,l)\mapsto(-p,-l)$.
Thus none of the complex-momentum formulas below uses literal
complex conjugation of its variable. Let $s$ be the coarse divisor, $b_a$ the
coarse bulb and $b_{l,a}$ the fine bulb. Write $q_0=p/L$ and
$u_a=\operatorname{tail}_a$, including the current noise. Splitting
the covariance recursion into its principal alias and the other
aliases gives
\begin{equation}\label{eq:transport-bulb-split}
 b_a=|m_0\operatorname{rail}_L(q_{0,a})|^2b_{0,a}+su_a.
\end{equation}
For $p\ne0$ this follows by multiplying
$D_a=|m_0\operatorname{rail}_L(q_{0,a})|^2b_{0,a}/s+u_a$
by $s$. Both sides of \eqref{eq:transport-bulb-split} are
continuous, so it holds at zero as well.

For the one-form column put $B_l=D_lF_{1,l}^*D^{-1}$ and
define $r_{l,a}$ by $B_lx=\operatorname{diag}(r_{l,a})x_l$.
Entrywise,
\[
 r_{l,a}=\overline{m_l}(D_l)_{aa}
        |\operatorname{rail}_L(q_{l,a})|^2/D_{aa}.
\]
Using \eqref{eq:transport-rotor} in \eqref{eq:transport-A} gives
\[
 \mathcal A_l=B_l+
                 \operatorname{diag}(e_l-r_{l,a})x_l\pi.
\]
Now $D^{-1}=s\operatorname{diag}(b_a^{-1})$ and
$\pi=s^{-1}\sigma$, where
$\sigma=(\operatorname{bend}_c)^{-1}
 x^*\operatorname{diag}(b_a^{-1})$ is analytic.
For a nonprincipal alias, $D_l$ has the regular port formula;
thus $e_l-r_{l,a}=s g_{l,a}$ with
\[
 g_{l,a}=\overline{m_l}
 \left(\frac{\operatorname{tr}D_l}{sa}
             -\frac{(D_l)_{aa}|\operatorname{rail}_L(q_{l,a})|^2}{b_a}
 \right).
\]
For the principal alias let $h$ be the sum of the nonprincipal
terms in $a$. Then $sa=|m_0|^2\sum_a b_{0,a}+sh$, and
\eqref{eq:transport-bulb-split} instead gives the regular formula
\[
       g_{0,a}=m_0^{-1}\left(\frac{u_a}{b_a}-\frac{h}{sa}\right).
\]
These are exactly the principal and nonprincipal gain formulas.
The denominators $b_a$, $sa$, and $m_0$ have the common analytic
lower bounds established in the component estimates. Consequently
\[
                \mathcal A_l=B_l+
                              \operatorname{diag}(g_{l,a})x_l\sigma
\]
is the already estimated analytic column, including at the origin.

For the two-form column, at a nonprincipal alias write
$\lambda_l=|m_l|^2\operatorname{bend}_c/\operatorname{portR}_l$.
Then $y_l=s^2\lambda_l$ and
\[
 s^2(I-C)=s^2I-sHd_1\operatorname{diag}(b_a)d_1^*.
\]
Substitution in \eqref{eq:transport-T} is exactly the regular
nonprincipal wire formula. At the principal alias set
$J=F_{2,0}^{-1}$, $U=\operatorname{diag}(u_a)$,
$M=\operatorname{diag}(b_a)$, and $z_a=u_a/b_a$.
The inverse exists on the common cube. The exterior chain rule
and the principal covariance split give
\[
 \delta_1=Jd_1F_{1,0},\qquad
 F_{1,0}D_0F_{1,0}^*=D-U.
\]
The cancellation needed for the apparent remaining pole is
\begin{equation}\label{eq:transport-brace-cancellation}
 s\mathcal B=H_0-y_0F_{2,0}^*HF_{2,0},\qquad
 \mathcal B=H_0\left(\operatorname{diag}_{\{a,b\}}(z_a+z_b)
                                  -s\bigwedge^2\operatorname{diag}(z_a)\right).
\end{equation}
To verify it, put
$Z_a=|m_0\operatorname{rail}_L(q_{0,a})|^2b_{0,a}/b_a
=1-sz_a$ by \eqref{eq:transport-bulb-split}.
The pulse formula shows
$y_0F_{2,0}^*HF_{2,0}=H_0\bigwedge^2\operatorname{diag}(Z_a)$:
on a face $\{a,b\}$ both sides have the same product of its two
diagonal entries and the same scalar bend ratio.
Finally $1-(1-sz_a)(1-sz_b)=s(z_a+z_b-sz_az_b)$ proves
\eqref{eq:transport-brace-cancellation} entrywise.

Using $M=sD$, the principal expression therefore becomes
\begin{equation}\label{eq:transport-principal-wire}
 \mathcal T_0=y_0F_{2,0}^*
               +(\mathcal B Jd_1M-H_0Jd_1U)d_1^*.
\end{equation}
Indeed its second term expands to
$H_0Jd_1(D-U)d_1^*-y_0F_{2,0}^*Hd_1Dd_1^*$,
which is \eqref{eq:transport-T} after the two chain identities.
All factors in \eqref{eq:transport-principal-wire} have the
analytic bounds already proved. This verifies agreement with the
principal wire as well as the nonprincipal one.

At $p=0$ the principal analytic columns have values
$\mathcal A_0=I$ and $\mathcal T_0=I$: the bulbs, bends,
rails, and principal filter have value one and the exterior
symbols vanish. For $l\ne0$ the patch mass filter vanishes at
$q_l(0)$. Indeed some nonzero alias coordinate has a nontrivial
$L$th root of unity, and the geometric-series rail vanishes
there. The corresponding plank is a nonzero phase times that rail, and
the patch mass is a positive power of the product of the four planks, so
it vanishes. The nonprincipal one-form column is zero because both
its leading term and gain contain that filter. The two-form
column is zero because its first term contains $d_1(0)^*=0$
and its other term contains the mass filter. Thus
\eqref{eq:transport-punctured-three} and
\eqref{eq:transport-alias-sum} also hold at $p=0$.

\subsection{Returning to physical counting measure}
The installed one-form and two-form columns at alias $l$ are
$L^{-1}\mathcal A_l$ and $L^{-2}\mathcal T_l$.
The block mean $Q$ is the alias row with columns $LF_{1,l}$.
In physical counting measure the adjoint alias column of $Q$
is $L^{-3}F_{1,l}^*$: the factor $L$ is multiplied by the
coarse-to-fine volume ratio $L^{-4}$. Since
$\delta_r=Ld_r(q_l)$, the second and third identities in
\eqref{eq:transport-punctured-three} become respectively
\[
 d_{1,f}^*T=Q^*d_{1,c}^*,\qquad
                 TH_cd_{1,c}=H_fd_{1,f}A_0.
\]
The alias sum becomes $QA_0=I-k^{-1}d_{1,c}^*H_cd_{1,c}$.
Fourier diagonalization on each finite torus now gives the real
matrix identities
\begin{equation}\label{eq:transport-installed-identities}
 T^{\mathsf T}d_{1,f}=d_{1,c}Q,\qquad
 TH_cd_{1,c}=H_fd_{1,f}A_0,\qquad
 QA_0=I-k^{-1}d_{1,c}^{\mathsf T}H_cd_{1,c}.
\end{equation}
This argument includes every finite period, including the one-site
period, and uses the Fourier coefficients of the original analytic
columns. No new choice of transport has been made.
